\apendice{Sostenibilización curricular}
 
\section{Introducción}
 
El presente apéndice analiza el Trabajo de Fin de Grado desde la perspectiva de
la sostenibilidad, tomando como referencia los \textbf{Objetivos de
Desarrollo Sostenible (ODS)} de la Agenda 2030 aprobada por la Organización de
las Naciones Unidas~\cite{onu:agenda2030}, en línea con las directrices de
sostenibilización curricular promovidas por CRUE Universidades
Españolas~\cite{crue:sostenibilidad}. Dado que el sistema HIVES se inscribe en el
ámbito agroalimentario y se vincula directamente con el sector apícola, el
análisis de su contribución a los objetivos resulta especialmente natural y
verificable, lo que permite una valoración fundamentada y evitando
declaraciones genéricas. En las secciones siguientes se identifican los objetivos
sobre los que el proyecto ejerce una contribución directa, se detallan aquellos
en los que influye y se examina, sin omitir sus límites, el impacto
del propio proceso de desarrollo.
 
\section{Contribución a los Objetivos de Desarrollo Sostenible}
 
La Agenda 2030 articula su llamamiento a la acción en torno a 17 Objetivos de
Desarrollo Sostenible interrelacionados, que integran de manera equilibrada las
dimensiones social, económica y ambiental del desarrollo~\cite{onu:agenda2030}.
La Figura \ref{fig:img/ods_agenda2030} recoge el conjunto completo de objetivos que conforman dicho
marco de referencia.
 
\imagentamano{img/ods_agenda2030}{Los 17 Objetivos de Desarrollo Sostenible de la Agenda 2030 de Naciones Unidas. Fuente: \cite{avv_ods}}{0.7}
 
El sistema HIVES no tiene como objetivo el influir sobre todos
los objetivos, planteamiento que transicionaría el análisis hacia irrelevante.
Por el contrario, se identifican \textbf{cuatro objetivos} sobre los que el
proyecto ejerce una contribución directa y demostrable, así como un conjunto
adicional de objetivos sobre los que incide de manera transversal.
 
\subsection{Objetivos abordados de forma directa}
 
La aportación  del proyecto se concentra en cuatro objetivos estrechamente
ligados a la naturaleza agroalimentaria y tecnológica de la solución, recogidos
en la Figura \ref{fig:img/ods_hives} y desarrollados a continuación en relación con sus metas
específicas.
\imagentamano{img/ods_hives}{Objetivos de Desarrollo Sostenible sobre los que el sistema HIVES ejerce una contribución directa. Fuente: Elaboración Propia.}{0.85}
 
\textbf{ODS 12 — Producción y consumo responsables.}
Constituye el objetivo más directamente vinculado al proyecto. Su
\textbf{meta 12.8} persigue garantizar que las personas dispongan de la
información pertinente para adoptar decisiones de consumo responsables. El fraude
alimentario y el etiquetado engañoso de la miel rompen precisamente esa cadena
de información, al desligar el precio y la etiqueta del origen botánico real del
producto. Al automatizar la verificación de dicho origen, HIVES refuerza la
\textbf{trazabilidad y la autenticidad} del producto, restituyendo al consumidor
una información veraz y reforzando la confianza en el mercado.
La solución incide igualmente en la meta 12.6, en tanto que dota a los agentes
productores de una herramienta para acreditar la calidad y la integridad de sus
productos.
 
\textbf{ODS 15 — Vida de ecosistemas terrestres.}
Su \textbf{meta 15.5} aborda la adopción de medidas para reducir la degradación
de los hábitats naturales y frenar la pérdida de biodiversidad. La conexión del
proyecto con este objetivo es de naturaleza económica e indirecta, pero sólida:
la viabilidad comercial de la apicultura constituye un incentivo decisivo para el
mantenimiento de las colmenas y, con ellas, de las poblaciones de abejas, agentes
polinizadores esenciales para la reproducción de numerosas especies vegetales y
para el sostenimiento de los ecosistemas. Al proteger el valor diferencial de las
mieles monoflorales frente a la adulteración, HIVES contribuye a preservar la
rentabilidad de la actividad de los apicultores.
 
\textbf{ODS 9 — Industria, innovación e infraestructura.}
Su \textbf{meta 9.5} promueve el aumento de la investigación científica y la
mejora de la capacidad tecnológica. El proyecto materializa una innovación
aplicada que \textbf{democratiza} el acceso a una técnica de análisis avanzada:
mediante la fusión de espectrometría óptica de bajo coste e inteligencia
artificial, traslada una capacidad analítica antes reservada a laboratorios
especializados a una solución autónoma de escritorio. Esta reducción de la
barrera tecnológica es, en sí misma, una contribución al desarrollo de
infraestructuras de calidad, fiables y asequibles.
 
\textbf{ODS 2 — Hambre cero.}
Más allá de la mera disponibilidad de alimentos, su \textbf{meta 2.4} incide en
la sostenibilidad y la fiabilidad de los sistemas de producción agroalimentaria.
La garantía de autenticidad y calidad de un producto de alto valor nutricional
como la miel se alinea con la promoción de cadenas alimentarias transparentes y
sostenibles, en las que la información del producto se corresponde con sus
características reales.
 
\subsection{Objetivos abordados de forma transversal}
 
Junto a los objetivos anteriores, el proyecto presenta vínculos secundarios con
otros ODS que, sin constituir su eje central, se ven favorecidos por su
desarrollo. Así, la defensa de la competencia leal y de los medios de vida de los
productores honestos frente a las prácticas adulteradoras conecta con el
\textbf{ODS 8} (trabajo decente y crecimiento económico); la detección de
adulteraciones que pudieran comprometer las propiedades del producto destinado al
consumo enlaza con el \textbf{ODS 3} (salud y bienestar); y el propio Trabajo de
Fin de Grado, en tanto que proceso formativo, ejercita competencias transversales
asociadas al \textbf{ODS 4} (educación de calidad).
 
La Tabla~\ref{tabla:ods-hives} consolida la totalidad de las contribuciones,
asociando cada objetivo con la evidencia concreta que la sustenta dentro del
proyecto.
 
\tablaSmall{Correspondencia entre los Objetivos de Desarrollo Sostenible y las aportaciones del proyecto HIVES.}
{p{1.6cm} p{3.4cm} p{8cm}}
{ods-hives}
{\textbf{ODS} & \textbf{Objetivo} & \textbf{Contribución del proyecto HIVES} \\}
{
12 & Producción y consumo responsables & Verificación automática de la autenticidad y el origen botánico; refuerzo de la trazabilidad y lucha contra el fraude alimentario (metas 12.8 y 12.6). \\
15 & Vida de ecosistemas terrestres    & Protección del valor económico de la apicultura como incentivo para la conservación de las abejas polinizadoras y la biodiversidad (meta 15.5). \\
9  & Industria, innovación e infraestructura & Innovación que democratiza una técnica analítica avanzada mediante hardware de bajo coste e inteligencia artificial (meta 9.5). \\
2  & Hambre cero                        & Garantía de calidad y autenticidad en un producto agroalimentario de alto valor, en favor de sistemas de producción fiables (meta 2.4). \\
8  & Trabajo decente y crecimiento económico & Defensa de la competencia leal y de los medios de vida de los productores honestos frente a las prácticas adulteradoras. \\
3  & Salud y bienestar                  & Detección de adulteraciones que podrían comprometer las propiedades del producto destinado al consumo. \\
4  & Educación de calidad               & El propio TFG, en tanto que proceso formativo, ejercita competencias transversales en sostenibilidad. \\
}
 
\section{Análisis del impacto del proyecto por dimensiones}
 
Con el fin de complementar el análisis anterior y evitar una valoración
exclusivamente propositiva, se examina el proyecto conforme a las tres dimensiones
clásicas de la sostenibilidad, contemplando tanto sus aportaciones como sus costes
reales. La Figura \ref{fig:img/sostenibilidad_dimensiones} ilustra dicha intersección.
 
\imagentamano{img/sostenibilidad_dimensiones}{Las tres dimensiones de la sostenibilidad y su intersección con el proyecto HIVES. Fuente: Elaboración Propia.}{0.6}
 
\textbf{Dimensión ambiental.}
El balance ambiental es marcadamente favorable. El método de análisis es
\textbf{no destructivo} y \textbf{libre de reactivos}, lo que elimina el consumo
de muestras y la generación de residuos químicos propios del procedimiento
melisopalinológico tradicional. En aras del rigor, debe reconocerse que la fase
de entrenamiento de los modelos conllevó un coste computacional no despreciable
derivado del uso de Google Colab Pro; dicho impacto queda acotado por
tratarse de una tarea puntual y no recurrente, y por desplegarse en producción un
modelo ligero que infiere localmente sobre CPU sin dependencia de la nube.
 
\textbf{Dimensión social.}
La principal aportación es la \textbf{democratización} del acceso a una capacidad
analítica antes ligada a instrumental avanzado y a personal especializado en
botánica. Al reducir la barrera técnica y económica, el sistema habilita a
pequeños productores y cooperativas para verificar la autenticidad de su
producción, contribuyendo a un mercado más transparente y equitativo, a la vez
que protege al consumidor frente al fraude.
 
\textbf{Dimensión económica.}
La viabilidad económica, detallada en el Apéndice~A, respalda la sostenibilidad a
largo plazo de la propuesta. Frente al elevado coste por muestra y a los tiempos
de días o semanas del análisis de referencia, el sistema ofrece un cribado de
coste marginal muy reducido y resultados en cuestión de minutos, con un coste
total de desarrollo estimado en 3.247,68 € y una arquitectura plenamente
\textbf{escalable}.
 
\section{Reflexión ética y profesional}
 
El desarrollo de una herramienta de verificación de la autenticidad de un
producto alimentario conlleva una responsabilidad deontológica que el proyecto
asume desde tres principios. El primero es la \textbf{honestidad metodológica}:
la documentación expone con transparencia las limitaciones del conjunto de datos,
el recurso al aumento sintético de muestras y las métricas reales alcanzadas, sin
presentar el prototipo como un sustituto con validez legal del método normalizado
de referencia. El segundo es la \textbf{minimización y custodia de datos}: la
información se persiste de forma local, sin transferencias a terceros, en
coherencia con la viabilidad legal analizada en el Apéndice~A. El tercero es la
\textbf{orientación al bien común}: la finalidad última de la herramienta es
proteger tanto al consumidor como al productor honesto, situando la tecnología al
servicio de la confianza en la cadena agroalimentaria.
 
\section{Conclusión}
 
El análisis evidencia que HIVES no constituye un desarrollo técnicamente neutro,
sino una solución con un impacto positivo articulado y verificable sobre la
sostenibilidad. Su contribución directa a los Objetivos de Desarrollo Sostenible
2, 9, 12 y 15, junto con su incidencia transversal sobre otros objetivos de la
Agenda 2030, inscribe el proyecto de pleno en el marco de la sostenibilización
curricular, integrando la responsabilidad ética y la conciencia de la huella
propia como componentes inseparables de la práctica profesional del ingeniero.